%=======================================================================
%  CHAPTER 2 — PROBLEM SOLVING : SOLVED PROBLEMS
%=======================================================================
\section{Problem Solving \textcolor{sub}{\large -- Solved Problems}}

{\color{sub}\footnotesize Every worked problem here opens with the exam question
\mk{exactly as the paper prints it}, typos and all, so it can be attempted
from this document alone. Methods follow \emph{Insights on Artificial Intelligence};
where that book is wrong or skips a step, the notes say so.
\gap$\cdot$\gap Every answer in this file is machine-checked by \code{verify.py}.}

\lead{The five things everything here is built on:}
\begin{itemize}
  \item \textbf{1. A crypt-arithmetic column is an equation.}
        $(\text{column sum}) + C_{\text{in}} = 10\,C_{\text{out}} + (\text{result digit})$.
        Every deduction in \S2.N.1 comes from one of these.
  \item \textbf{2. A carry is 0 or 1.} Adding two digits plus a carry gives at most
        $9+9+1=19$. With \emph{three} addends it can reach $9+9+9+1=28$, so a carry of
        \textbf{2} is possible in TEN $+$ TEN $+$ FORTY. Nothing else in these papers has
        three addends.
  \item \textbf{3. If the answer is one digit longer than the addends, that leading digit
        is 1} $-$ two 4-digit numbers cannot reach 20000. This single line gives
        $M = 1$ in SEND $+$ MORE, $G = 1$ in BASE $+$ BALL, $P = 1$ in LOGIC $+$ LOGIC and
        $D = 1$ in CROSS $+$ ROADS, for free.
  \item \textbf{4. $2X$ ends in $X$ only when $X = 0$.} ($X=5$ gives 10, whose last digit
        is 0, not 5.) This cracks TEN $+$ TEN $+$ FORTY and LOVE $+$ LOVE.
  \item \textbf{5. A jug state is the ordered pair of contents.} Every water-jug operator
        maps $(x,y)$ to another $(x,y)$; there is nothing else to track.
\end{itemize}

\hr

\lead{\mk{Read this before the crypt-arithmetic problems}}
\begin{itemize}
  \item \mk{Seven of the thirteen puzzles set in these papers have more than
        one valid answer.} ONE $+$ ONE $+$ TWO $=$ FOUR has \textbf{107}.
        RIGHT $+$ RIGHT $=$ WRONG, WRONG $+$ WRONG $=$ RIGHT and ONE $+$ ONE $=$ TWO pin
        \textbf{no letter at all}.
  \item So the marks are in the \textbf{column reasoning}, not the digits. Write $X$, $D$,
        $C$ and the carries, then argue. Where the constraints run out, say
        \emph{``the constraints leave $N$ free; take $N = 3$''} and carry on. That sentence
        is worth more than the answer.
  \item Uniqueness of each puzzle is stated below it, counted by exhaustive search in
        \code{verify.py} over every letter-to-digit assignment.
\end{itemize}

\hr

%=======================================================================
\T{2.N.1 Crypt-arithmetic: the six puzzles with a unique answer}

%-----------------------------------------------------------------------
\creamq{SEND $+$ MORE $=$ MONEY}{\m{1+6}\m{8}}
{\color{sub}\footnotesize \tF{5} \yr{\textbf{72 Ash} \m{1+6}, 76 Ba \m{8},
\textit{74 Ma} \m{8}, \textbf{\textit{74 Bh}} \m{8}, \textbf{\textit{72 Ash}} \m{8}}
\gap$\cdot$\gap \mk{unique solution} \gap$\cdot$\gap the classic, and the one
\emph{Insights} works in full}

\begin{asked}{\textbf{72 Ash} \textperiodcentered{} Q2 \textperiodcentered{} \m{1+6}}
Solve the following crypto-arithmetic problem through constraint satisfaction, showing all
steps. Different letters denote different integers and identical letters denote the same
integer.
\begin{center}
\begin{tabular}{r@{\hspace{2pt}}c@{\hspace{4pt}}c@{\hspace{4pt}}c@{\hspace{4pt}}c@{\hspace{4pt}}c}
   &   & S & E & N & D \\
 $+$ &   & M & O & R & E \\
\hline
   & M & O & N & E & Y \\
\end{tabular}
\end{center}
\end{asked}

\sbsr{0.46}{%
\lead{Formulation}
\begin{itemize}
  \item $X = \{S, E, N, D, M, O, R, Y\}$ \quad (8 letters)
  \item $D_{\text{om}} = \{0,1,2,3,4,5,6,7,8,9\}$
  \item $C$:
  \begin{itemize}
    \item $S \neq 0$, $M \neq 0$ (leading letters)
    \item all eight letters take different digits
    \item $D + E = 10C_1 + Y$
    \item $N + R + C_1 = 10C_2 + E$
    \item $E + O + C_2 = 10C_3 + N$
    \item $S + M + C_3 = 10C_4 + O$
    \item $C_4 = M$
  \end{itemize}
\end{itemize}
\vspace{2pt}
\begin{center}
\begin{tabular}{r@{\hspace{3pt}}c@{\hspace{5pt}}c@{\hspace{5pt}}c@{\hspace{5pt}}c@{\hspace{5pt}}c}
 & $C_4$ & $C_3$ & $C_2$ & $C_1$ & \\
 &   & S & E & N & D \\
$+$ &   & M & O & R & E \\
\hline
 & M & O & N & E & Y \\
\end{tabular}
\end{center}}{%
\lead{The deduction chain}
\begin{itemize}
  \item \textbf{$M = 1$.} $C_4 = M$, and a carry is at most 1. $M \neq 0$ because it
        leads MONEY, so $M = 1$.
  \item \textbf{$O = 0$.} Column 4 is $S + 1 + C_3 = 10 \cdot 1 + O$. The largest
        $S + 1 + C_3$ can be is $9+1+1 = 11$, so $O$ is 0 or 1; 1 is taken by $M$, so
        $O = 0$.
  \item \textbf{$S = 9$.} With $O = 0$, column 4 reads $S + 1 + C_3 = 10$. Column 3 is
        $E + 0 + C_2 = 10C_3 + N$, and $E + C_2 \leq 10$, so $C_3 = 0$ unless $E+C_2=10$,
        which would force $N = 0 = O$. So $C_3 = 0$ and $S = 9$.
  \item \textbf{$N = E + 1$ and $C_2 = 1$.} Column 3 now reads $E + C_2 = N$. If
        $C_2$ were 0 then $E = N$, not allowed. So $C_2 = 1$ and $N = E + 1$.
  \item \textbf{$R = 8$.} Column 2 is $N + R + C_1 = 10 + E$. Substitute $N = E+1$:
        $E + 1 + R + C_1 = E + 10$, so $R + C_1 = 9$. $R = 9$ is taken by $S$, so
        $C_1 = 1$ and $R = 8$.
  \item \textbf{$D = 7$.} Column 1 is $D + E = 10 + Y$ (since $C_1 = 1$), so
        $D + E \geq 12$. With 9 and 8 gone, the only pairs are $5+7$ and $6+7$, so one of
        $D, E$ is 7. If $E = 7$ then $N = 8 = R$, so \textbf{$D = 7$}, and $E$ is 5 or 6.
  \item \textbf{$E = 5$, $N = 6$, $Y = 2$.} $E = 6$ would give $N = 7 = D$. So $E = 5$,
        $N = 6$, and $Y = D + E - 10 = 2$.
\end{itemize}}

\penalty0\vspace{2pt}
{\color{acc}\bfseries Answer:} $S{=}9,\ E{=}5,\ N{=}6,\ D{=}7,\ M{=}1,\ O{=}0,\ R{=}8,\ Y{=}2$
\quad$\rightarrow$\quad $9567 + 1085 = 10652$. \quad
{\color{sub}\footnotesize Unique: exactly one assignment satisfies the constraints.}

%-----------------------------------------------------------------------
\creamq{BASE $+$ BALL $=$ GAMES}{\m{1+6}\m{2+6}\m{3+5}\m{8}}
{\color{sub}\footnotesize \tF{7} \yr{75 Ba \m{2+6}, 82 Ka \m{2+6}, \texttt{82 Ba} \m{2+6},
79 Ash \m{2+6}, \textbf{\texttt{82 Bh}} \m{3+5}, 79 Jth \m{1+6},
\textbf{\textit{73 Bh}} \m{8}} \gap$\cdot$\gap
\mk{the most-asked puzzle in the subject} \gap$\cdot$\gap unique solution}

\begin{asked}{79 Ash \textperiodcentered{} \m{2+6} \gap$\cdot$\gap also 75 Ba, 82 Ka, \texttt{82 Ba}, \textbf{\texttt{82 Bh}}, 79 Jth, \textbf{\textit{73 Bh}}}
What do you understand by Constraint Satisfaction Problem? Solve the following
crypto-arithmetic problem through constraint satisfaction, showing all steps. Different
letters denote different integers and identical letters denote the same integer.
\begin{center}
\begin{tabular}{r@{\hspace{2pt}}c@{\hspace{4pt}}c@{\hspace{4pt}}c@{\hspace{4pt}}c@{\hspace{4pt}}c}
   &   & B & A & S & E \\
 $+$ &   & B & A & L & L \\
\hline
   & G & A & M & E & S \\
\end{tabular}
\end{center}
\end{asked}

\sbsr{0.44}{%
\lead{Formulation}
\begin{itemize}
  \item $X = \{B, A, S, E, L, G, M\}$ \quad (7 letters)
  \item $D_{\text{om}} = \{0,\dots,9\}$
  \item $C$:
  \begin{itemize}
    \item $B \neq 0$, $G \neq 0$
    \item all seven letters differ
    \item $E + L = 10C_1 + S$
    \item $S + L + C_1 = 10C_2 + E$
    \item $A + A + C_2 = 10C_3 + M$
    \item $B + B + C_3 = 10C_4 + A$
    \item $C_4 = G$
  \end{itemize}
\end{itemize}
\vspace{2pt}
\begin{center}
\begin{tabular}{r@{\hspace{3pt}}c@{\hspace{5pt}}c@{\hspace{5pt}}c@{\hspace{5pt}}c@{\hspace{5pt}}c}
 & $C_4$ & $C_3$ & $C_2$ & $C_1$ & \\
 &   & B & A & S & E \\
$+$ &   & B & A & L & L \\
\hline
 & G & A & M & E & S \\
\end{tabular}
\end{center}}{%
\lead{The deduction chain}
\begin{itemize}
  \item \textbf{$G = 1$.} Two 4-digit numbers cannot reach 20000, so the leading carry
        $C_4 = G = 1$.
  \item \textbf{$L = 5$, $C_1 = 0$, $C_2 = 1$.} \mk{This is the clever step.}
        Add the column-1 and column-2 equations:
        \begin{center}$(E + L) + (S + L + C_1) = (10C_1 + S) + (10C_2 + E)$\end{center}
        Everything cancels except $2L = 9C_1 + 10C_2$. With $2L \leq 18$ the only
        possibilities are $(C_1{=}0, C_2{=}0, L{=}0)$ and $(C_1{=}0, C_2{=}1, L{=}5)$;
        $2L = 9$ and $2L = 19$ are odd, so no other case survives.
        $L = 0$ makes column 1 read $E = S$, which the AllDiff constraint forbids.
        So \textbf{$L = 5$}.
  \item \textbf{$S = E + 5$.} Column 1 with $L=5$, $C_1=0$ is $E + 5 = S$. Hence
        $E \leq 4$, and $E \neq 0$ (else $S = 5 = L$), so $E \in \{2,3,4\}$ and
        $S \in \{7,8,9\}$.
  \item \textbf{Columns 3 and 4.} $2A + 1 = 10C_3 + M$ and $2B + C_3 = A + 10$.
  \item \textbf{Test the three cases for $E$:}
  \begin{itemize}
    \item $E = 2, S = 7$: $C_3{=}0$ forces $M = 9, A = 4$, then $2B = 14$, $B = 7 = S$.
          Clash. $C_3{=}1$ gives no legal $B$. \xmark
    \item $E = 4, S = 9$: every branch needs $2B$ odd, or repeats $G$. \xmark
    \item $E = 3, S = 8$: $C_3 = 0$, $M = 2A+1$ odd and unused $\rightarrow M = 9$,
          $A = 4$; then $2B + 0 = 14$, so $\mathbf{B = 7}$. \checkmark
  \end{itemize}
\end{itemize}}

\penalty0\vspace{2pt}
{\color{acc}\bfseries Answer:} $B{=}7,\ A{=}4,\ S{=}8,\ E{=}3,\ L{=}5,\ G{=}1,\ M{=}9$
\quad$\rightarrow$\quad $7483 + 7455 = 14938$. \quad
{\color{sub}\footnotesize Unique.}

%-----------------------------------------------------------------------
\creamq{TEN $+$ TEN $+$ FORTY $=$ SIXTY}{\m{2+6}}
{\color{sub}\footnotesize \tP{3} \yr{\textbf{\texttt{80 Bh}} \m{2+6}, 78 Ba \m{2+6},
78 Po \m{2+6} (printed FORTY $+$ TEN $+$ TEN, same puzzle)} \gap$\cdot$\gap
\mk{the only puzzle with three addends}, so here a carry can be 2 \gap$\cdot$\gap
unique solution}

\begin{asked}{\textbf{\texttt{80 Bh}} \textperiodcentered{} \m{2+6} \gap$\cdot$\gap 78 Ba; 78 Po prints the addends as FORTY $+$ TEN $+$ TEN}
What do you understand by Constraint Satisfaction Problem? Solve the following
crypto-arithmetic problem through constraint satisfaction, showing all steps.
\begin{center}
\begin{tabular}{r@{\hspace{2pt}}c@{\hspace{4pt}}c@{\hspace{4pt}}c@{\hspace{4pt}}c@{\hspace{4pt}}c}
   &   &   & T & E & N \\
 $+$ &   &   & T & E & N \\
 $+$ & F & O & R & T & Y \\
\hline
   & S & I & X & T & Y \\
\end{tabular}
\end{center}
\end{asked}

\sbsr{0.44}{%
\lead{Formulation}
\begin{itemize}
  \item $X = \{T,E,N,F,O,R,Y,S,I,X\}$ \quad
        \mk{all 10 letters}, so every digit 0--9 is used exactly once.
  \item $C$: $T, F, S \neq 0$; all ten differ; and
  \begin{itemize}
    \item $N + N + Y = 10C_1 + Y$
    \item $E + E + T + C_1 = 10C_2 + T$
    \item $T + T + R + C_2 = 10C_3 + X$
    \item $O + C_3 = 10C_4 + I$
    \item $F + C_4 = S$
  \end{itemize}
  \item \mk{Note $C_3$ may be 2}: column 3 can reach $9+9+9+1 = 28$.
\end{itemize}
\vspace{2pt}
\begin{center}
\begin{tabular}{r@{\hspace{3pt}}c@{\hspace{5pt}}c@{\hspace{5pt}}c@{\hspace{5pt}}c@{\hspace{5pt}}c}
 & $C_4$ & $C_3$ & $C_2$ & $C_1$ & \\
 &   &   & T & E & N \\
$+$ &   &   & T & E & N \\
$+$ & F & O & R & T & Y \\
\hline
 & S & I & X & T & Y \\
\end{tabular}
\end{center}}{%
\lead{The deduction chain}
\begin{itemize}
  \item \textbf{$N = 0$ and $C_1 = 0$.} Column 1 is $2N + Y = 10C_1 + Y$, so $2N = 10C_1$.
        Hence $N = 0$ (with $C_1 = 0$) or $N = 5$ (with $C_1 = 1$).
  \item \textbf{$E = 5$ and $C_2 = 1$.} Column 2 is $2E + T + C_1 = 10C_2 + T$, so
        $2E + C_1 = 10C_2$ by the same cancellation. If $N = 5$ then $C_1 = 1$ and
        $2E = 10C_2 - 1$ is odd: impossible. So \textbf{$N = 0$, $C_1 = 0$}, and then
        $2E = 10C_2$ gives $E = 5$, $C_2 = 1$ ($E = 0$ is taken by $N$).
  \item \textbf{$C_4 = 1$ and $S = F + 1$.} $F$ and $S$ must differ, so the last column
        $F + C_4 = S$ needs $C_4 = 1$: $S$ and $F$ are \textbf{consecutive}.
  \item \textbf{$O = 9$, $C_3 = 2$, $I = 1$.} Column 4 is $O + C_3 = 10 + I$ (since
        $C_4 = 1$), so $O + C_3 \geq 10$ and $O$ is large. $O = 8$ with $C_3 = 2$ gives
        $I = 0 = N$; $O = 9$ with $C_3 = 1$ gives $I = 0 = N$. So \textbf{$O = 9$,
        $C_3 = 2$, $I = 1$}.
  \item \textbf{$T = 8$, $R = 7$, $X = 4$.} Column 3 is $2T + R + 1 = 20 + X$. Digits
        $0,1,5,9$ are used, leaving $\{2,3,4,6,7,8\}$ for $T,R,X,F,S,Y$. $2T + R = 19 + X$
        needs $T$ large: $T = 8$ gives $R = X + 3$, and $(R,X) = (7,4)$ is the only pair
        left in the set.
  \item \textbf{$F = 2$, $S = 3$, $Y = 6$.} $\{2,3,6\}$ remain. $S = F + 1$ forces
        $F = 2$, $S = 3$, leaving $Y = 6$.
\end{itemize}}

\penalty0\vspace{2pt}
{\color{acc}\bfseries Answer:} $T{=}8,\ E{=}5,\ N{=}0,\ F{=}2,\ O{=}9,\ R{=}7,\ Y{=}6,\ S{=}3,\ I{=}1,\ X{=}4$
\quad$\rightarrow$\quad $850 + 850 + 29786 = 31486$. \quad
{\color{sub}\footnotesize Unique.}

{\color{sub}\footnotesize \mk{Where \emph{Insights} is loose here}: it writes
``Let us take $R = 7$, then $X = 4$'' as though it were a free choice, and produces $T = 8$
without saying where it came from. The set $\{2,3,4,6,7,8\}$ and the equation
$2T + R = 19 + X$ actually force all three; the step above supplies the missing argument.}

%-----------------------------------------------------------------------
\creamq{LOGIC $+$ LOGIC $=$ PROLOG}{\m{3+4}\m{3+5}\m{7}}
{\color{sub}\footnotesize \tP{3} \yr{\textbf{69 Bh} \m{7}, 72 Ma \m{3+4},
\textbf{79 Ch} \m{3+5}} \gap$\cdot$\gap unique solution}

\begin{asked}{\textbf{79 Ch} \textperiodcentered{} \m{3+5} \gap$\cdot$\gap also \textbf{69 Bh} \m{7}, 72 Ma \m{3+4}}
Solve the following crypto-arithmetic problem through constraint satisfaction, showing all
steps. Different letters denote different integers and identical letters denote the same
integer.
\begin{center}
\begin{tabular}{r@{\hspace{2pt}}c@{\hspace{4pt}}c@{\hspace{4pt}}c@{\hspace{4pt}}c@{\hspace{4pt}}c@{\hspace{4pt}}c}
   &   & L & O & G & I & C \\
 $+$ &   & L & O & G & I & C \\
\hline
   & P & R & O & L & O & G \\
\end{tabular}
\end{center}
\end{asked}

\sbsr{0.44}{%
\lead{Formulation}
\begin{itemize}
  \item $X = \{L,O,G,I,C,P,R\}$; $L, P \neq 0$; all seven differ.
  \item The sum is $2 \times \text{LOGIC} = \text{PROLOG}$, columns right to left:
  \begin{itemize}
    \item $2C = 10a_1 + G$
    \item $2I + a_1 = 10a_2 + O$
    \item $2G + a_2 = 10a_3 + L$
    \item $2O + a_3 = 10a_4 + O$
    \item $2L + a_4 = 10a_5 + R$
    \item $a_5 = P$
  \end{itemize}
\end{itemize}
\vspace{2pt}
\begin{center}
\begin{tabular}{r@{\hspace{3pt}}c@{\hspace{4pt}}c@{\hspace{4pt}}c@{\hspace{4pt}}c@{\hspace{4pt}}c@{\hspace{4pt}}c}
 & $a_5$ & $a_4$ & $a_3$ & $a_2$ & $a_1$ & \\
 &  & L & O & G & I & C \\
$+$ &  & L & O & G & I & C \\
\hline
 & P & R & O & L & O & G \\
\end{tabular}
\end{center}}{%
\lead{The deduction chain}
\begin{itemize}
  \item \textbf{$P = 1$.} $2 \times$ a 5-digit number is under 200000, so the leading
        carry $a_5 = P = 1$.
  \item \textbf{$O = 0$, $a_3 = a_4 = 0$.} Column 4 is $2O + a_3 = 10a_4 + O$, i.e.\
        $O + a_3 = 10a_4$. Either $a_4 = 0$, forcing $O = a_3 = 0$; or $a_4 = 1$ with
        $O = 9, a_3 = 1$. \mk{Take the first branch and check the second at
        the end} $-$ the second dies below.
  \item \textbf{$I = 5$, $a_1 = 0$, $a_2 = 1$.} Column 2 is $2I + a_1 = 10a_2 + O = 10a_2$.
        $a_2 = 0$ would give $I = 0 = O$, so $a_2 = 1$ and $2I + a_1 = 10$; $a_1$ must be
        even, so $a_1 = 0$ and $I = 5$.
  \item \textbf{$G = 2C$ and $L = 2G + 1$.} Column 1 with $a_1 = 0$ gives $G = 2C$.
        Column 3 with $a_3 = 0$ gives $L = 2G + a_2 = 2G + 1 = 4C + 1$.
  \item \textbf{$R = 8C - 8$.} Column 5 with $a_4 = 0$, $a_5 = 1$ gives $2L = 10 + R$, so
        $R = 2L - 10 = 8C - 8$.
  \item \textbf{$C = 2$.} $L = 4C + 1 \leq 9$ forces $C \leq 2$. $C = 0$ clashes with
        $O$; $C = 1$ gives $L = 5 = I$. So \textbf{$C = 2$}, hence $G = 4$, $L = 9$,
        $R = 8$.
  \item \textbf{The other branch dies.} With $O = 9, a_3 = 1$: column 2 forces $I = 4$,
        $a_1 = 1$; column 1 then gives $G = 2C - 10$ and column 3 gives $L = 2G - 10 =
        4C - 30$, needing $C \geq 8$. $C = 8$ makes $R = 2L + 1 - 10 < 0$ and $C = 9$
        clashes with $O$. \xmark
\end{itemize}}

\penalty0\vspace{2pt}
{\color{acc}\bfseries Answer:} $L{=}9,\ O{=}0,\ G{=}4,\ I{=}5,\ C{=}2,\ P{=}1,\ R{=}8$
\quad$\rightarrow$\quad $90452 + 90452 = 180904$. \quad
{\color{sub}\footnotesize Unique.}

%-----------------------------------------------------------------------
\creamq{EAT $+$ THAT $=$ APPLE}{\m{2+5}\m{2+6}}
{\color{sub}\footnotesize \tP{2} \yr{\texttt{80 Ba} \m{2+6}, \textbf{80 Ch} \m{2+5}}
\gap$\cdot$\gap \mk{the only puzzle whose addends differ in length}
\gap$\cdot$\gap unique solution}

\begin{asked}{\texttt{80 Ba} \textperiodcentered{} \m{2+6} \gap$\cdot$\gap also \textbf{80 Ch} \m{2+5}}
What do you understand by Constraint Satisfaction Problem? Solve the following
crypto-arithmetic problem through constraint satisfaction, showing all steps.
\begin{center}
\begin{tabular}{r@{\hspace{2pt}}c@{\hspace{4pt}}c@{\hspace{4pt}}c@{\hspace{4pt}}c@{\hspace{4pt}}c}
   &   &   & E & A & T \\
 $+$ &   & T & H & A & T \\
\hline
   & A & P & P & L & E \\
\end{tabular}
\end{center}
\end{asked}

\sbsr{0.44}{%
\lead{Formulation}
\begin{itemize}
  \item $X = \{E,A,T,H,P,L\}$; $E, T, A \neq 0$; all six differ.
  \item Columns right to left:
  \begin{itemize}
    \item $T + T = 10C_1 + E$
    \item $A + A + C_1 = 10C_2 + L$
    \item $E + H + C_2 = 10C_3 + P$
    \item $T + C_3 = 10C_4 + P$
    \item $C_4 = A$
  \end{itemize}
\end{itemize}
\vspace{2pt}
\begin{center}
\begin{tabular}{r@{\hspace{3pt}}c@{\hspace{5pt}}c@{\hspace{5pt}}c@{\hspace{5pt}}c@{\hspace{5pt}}c}
 & $C_4$ & $C_3$ & $C_2$ & $C_1$ & \\
 &  &  & E & A & T \\
$+$ &  & T & H & A & T \\
\hline
 & A & P & P & L & E \\
\end{tabular}
\end{center}}{%
\lead{The deduction chain}
\begin{itemize}
  \item \textbf{$A = 1$.} A 3-digit plus a 4-digit number is at most
        $999 + 9999 = 10998$, so the 5-digit answer begins with 1.
  \item \textbf{$T = 9$.} The sum is at least 10000, so
        $\text{THAT} \geq 10000 - 999 = 9001$, forcing $T = 9$.
  \item \textbf{$E = 8$, $C_1 = 1$.} Column 1 is $9 + 9 = 18$, so $E = 8$ with a carry
        of 1.
  \item \textbf{$L = 3$, $C_2 = 0$.} Column 2 is $1 + 1 + 1 = 3$, so $L = 3$, no carry.
  \item \textbf{$P = 0$, $C_3 = 1$.} Column 4 is $T + C_3 = 10C_4 + P$ with $C_4 = A = 1$,
        i.e.\ $9 + C_3 = 10 + P$, so $P = C_3 - 1$. A carry is 0 or 1, so $C_3 = 1$ and
        \textbf{$P = 0$}.
  \item \textbf{$H = 2$.} Column 3 is $E + H + C_2 = 10C_3 + P$, i.e.\
        $8 + H + 0 = 10 + 0$, so $H = 2$.
  \item \mk{Every letter fell out of the two length observations.} No
        case-splitting was needed at all, which is why this puzzle is worth doing first if
        you are short of time.
\end{itemize}}

\penalty0\vspace{2pt}
{\color{acc}\bfseries Answer:} $E{=}8,\ A{=}1,\ T{=}9,\ H{=}2,\ P{=}0,\ L{=}3$
\quad$\rightarrow$\quad $819 + 9219 = 10038$. \quad
{\color{sub}\footnotesize Unique.}

%-----------------------------------------------------------------------
\creamq{CROSS $+$ ROADS $=$ DANGER}{\m{2+6}}
{\color{sub}\footnotesize \tP{1} \yr{\textbf{\texttt{81 Bh}} \m{2+6}} \gap$\cdot$\gap
unique solution}

\begin{asked}{\textbf{\texttt{81 Bh}} \textperiodcentered{} \m{2+6}}
What do you understand by Constraint Satisfaction Problem? Discuss it, and list all the
constraints. Solve the following crypto-arithmetic problem through constraint
satisfaction, showing all steps.
\begin{center}
\begin{tabular}{r@{\hspace{2pt}}c@{\hspace{4pt}}c@{\hspace{4pt}}c@{\hspace{4pt}}c@{\hspace{4pt}}c@{\hspace{4pt}}c}
   &   & C & R & O & S & S \\
 $+$ &   & R & O & A & D & S \\
\hline
   & D & A & N & G & E & R \\
\end{tabular}
\end{center}
\end{asked}

\sbsr{0.44}{%
\lead{Formulation}
\begin{itemize}
  \item $X = \{C,R,O,S,A,D,N,G,E\}$ \quad (9 letters, so one digit goes unused)
  \item $C, R, D \neq 0$; all nine differ.
  \item Columns right to left:
  \begin{itemize}
    \item $S + S = 10a_1 + R$
    \item $S + D + a_1 = 10a_2 + E$
    \item $O + A + a_2 = 10a_3 + G$
    \item $R + O + a_3 = 10a_4 + N$
    \item $C + R + a_4 = 10a_5 + A$
    \item $a_5 = D$
  \end{itemize}
\end{itemize}}{%
\lead{The deduction chain}
\begin{itemize}
  \item \textbf{$D = 1$.} Two 5-digit numbers cannot reach 200000, so the leading carry
        $a_5 = D = 1$.
  \item \textbf{$R$ is even.} $R$ is the last digit of $2S$, and $2S$ is even.
  \item \textbf{$C + R + a_4 = A + 10$}, from the leftmost column with $a_5 = 1$. Since
        $C, R \leq 9$, this needs $C$ and $R$ both large.
  \item These three plus AllDiff leave a small search. Propagating them column by column
        (the constraint-propagation loop of \S2.5) terminates on exactly one assignment,
        below.
  \item \mk{The column check, which is the part that earns the marks:}
  \begin{itemize}
    \item col 1: $3 + 3 = 6 = R$, \quad $a_1 = 0$ \checkmark
    \item col 2: $3 + 1 + 0 = 4 = E$, \quad $a_2 = 0$ \checkmark
    \item col 3: $2 + 5 + 0 = 7 = G$, \quad $a_3 = 0$ \checkmark
    \item col 4: $6 + 2 + 0 = 8 = N$, \quad $a_4 = 0$ \checkmark
    \item col 5: $9 + 6 + 0 = 15 = A + 10 = 5 + 10$, \quad $a_5 = 1 = D$ \checkmark
  \end{itemize}
  \item Digit 0 is the one left unused.
\end{itemize}}

\penalty0\vspace{2pt}
{\color{acc}\bfseries Answer:} $C{=}9,\ R{=}6,\ O{=}2,\ S{=}3,\ A{=}5,\ D{=}1,\ N{=}8,\ G{=}7,\ E{=}4$
\quad$\rightarrow$\quad $96233 + 62513 = 158746$. \quad
{\color{sub}\footnotesize Unique.}

%=======================================================================
\T{2.N.2 Crypt-arithmetic: the seven puzzles with more than one answer}

{\color{sub}\footnotesize These are set just as often as the unique ones, and a marker
cannot be expecting a particular set of digits. \mk{Show the forced letters,
name the free ones, choose values, and verify.} The counts below are from exhaustive
search in \code{verify.py}.}

\vspace{3pt}
\begin{tabularx}{\textwidth}{@{}L{3.75cm}L{1.4cm}L{2.9cm}X@{}}
\toprule
\hrow \thead{Puzzle} & \thead{Answers} & \thead{Forced by the constraints} & \thead{A valid assignment, and why the rest is free} \\
\midrule
ONE $+$ ONE $+$ TWO $=$ FOUR
  \newline {\footnotesize\color{sub}71 Ma, \textbf{73 Bh}, 80 Ash}
  & \mk{107} & $F = 1$ only
  & $O{=}6, N{=}4, E{=}2, T{=}3, W{=}8, F{=}1, U{=}7, R{=}0$: $642+642+386=1670$.
    Eight letters, but only the leading-carry constraint bites. \\
LOVE $+$ LOVE $=$ HATE
  \newline {\footnotesize\color{sub}\texttt{81 Ba}}
  & 42 & $E = 0$ only
  & $L{=}3, O{=}5, V{=}6, E{=}0, H{=}7, A{=}1, T{=}2$: $3560+3560=7120$.
    $E = 0$ because $2E$ must end in $E$ (rule 4). \\
WRONG $+$ WRONG $=$ RIGHT
  \newline {\footnotesize\color{sub}\textbf{70 Bh}, \textbf{77 Ch}}
  & 21 & \mk{nothing}
  & $W{=}3, R{=}7, O{=}0, N{=}8, G{=}1, I{=}4, H{=}6, T{=}2$: $37081+37081=74162$. \\
SWIM $+$ WEAR $=$ RELAX
  \newline {\footnotesize\color{sub}\textbf{74 Bh}, \textbf{81 Ch}}
  & 16 & $R = 1$, $I = 0$
  & $S{=}9, W{=}4, I{=}0, M{=}5, E{=}3, A{=}2, R{=}1, L{=}7, X{=}6$: $9405+4321=13726$.
    $R=1$ is the leading carry. \\
ONE $+$ ONE $=$ TWO
  \newline {\footnotesize\color{sub}\emph{Insights} \S2.3}
  & 16 & \mk{nothing}
  & $O{=}2, N{=}3, E{=}1, T{=}4, W{=}6$: $231+231=462$. This is the book's own worked
    example, and it presents its answer as though it were the only one. \\
RIGHT $+$ RIGHT $=$ WRONG
  \newline {\footnotesize\color{sub}69 Po}
  & 11 & \mk{nothing}
  & $R{=}4, I{=}2, G{=}0, H{=}6, T{=}5, W{=}8, O{=}1, N{=}3$: $42065+42065=84130$. \\
TWO $+$ TWO $=$ FOUR
  \newline {\footnotesize\color{sub}\textbf{78 Ch}}
  & 7 & $F = 1$ only
  & $T{=}7, W{=}3, O{=}4, F{=}1, U{=}6, R{=}8$: $734+734=1468$. \\
\bottomrule
\end{tabularx}

\penalty0\vspace{3pt}
\sbs{%
\lead{How to write one of these up for full marks}
\begin{itemize}
  \item \textbf{1.} Give $X$, $D$ and $C$ with the carry equations, exactly as for a unique
        puzzle. \mk{This is most of the mark scheme.}
  \item \textbf{2.} Derive whatever \emph{is} forced. There is always something: a leading
        carry, or a $2X$ ending in $X$.
  \item \textbf{3.} Say plainly that the constraints do not determine the rest, then pick
        values that satisfy every column and the AllDiff rule.
  \item \textbf{4.} \textbf{Verify by doing the addition.} One line, and it converts a
        guess into a demonstrated answer.
\end{itemize}}{%
\lead{Worked example of step 2: LOVE $+$ LOVE $=$ HATE}
\begin{itemize}
  \item Column 1: $E + E = 10C_1 + E \Rightarrow E = 10C_1$. A digit, so
        \textbf{$E = 0$, $C_1 = 0$}.
  \item Column 2: $2V + 0 = 10C_2 + T$, so $T$ is even.
  \item Column 3: $2O + C_2 = 10C_3 + A$.
  \item Column 4: $2L + C_3 = H$, and $H \leq 9$ forces $\mathbf{L \leq 4}$.
  \item That is everything the constraints give. $L \in \{1,2,3,4\}$ and the rest follows
        the choice. Take $L = 3, O = 5, V = 6$: then $T = 2$, $A = 1$, $H = 7$, and
        $3560 + 3560 = 7120$ \checkmark
\end{itemize}}

%=======================================================================
\T{2.N.2b The deduction chain for each of the five remaining puzzles}

{\color{sub}\footnotesize The table above is an index, not an answer. \mk{Nine
papers set one of these puzzles for six or seven marks}, and the marks are in the
\emph{chain}, not in the digits at the end of it $-$ a puzzle with 107 solutions cannot
be marked on the digits. LOVE $+$ LOVE $=$ HATE is chained just above; the other five
follow, each with its own paper, its own column equations and its own worked choice.
\gap$\cdot$\gap \textbf{Column $k$ is counted from the units end}, and $c_k$ is the carry
\emph{out} of column $k$. Every forced result below was confirmed against the full
solution set by \code{verify.py}.}

%-----------------------------------------------------------------------
\creamq{ONE $+$ ONE $+$ TWO $=$ FOUR}{\m{7}\m{2+6}}
{\color{sub}\footnotesize \tP{3} \yr{71 Ma \m{7}, \textbf{73 Bh} \m{7}, 80 Ash \m{2+6}}
\gap$\cdot$\gap \mk{107 solutions}, the most of any puzzle in the 41 papers
\gap$\cdot$\gap the only \emph{multi-answer} puzzle with three addends, so a carry here
can be \textbf{2}, not just 0 or 1}

\begin{asked}{\textbf{73 Bh} \textperiodcentered{} Q2 \textperiodcentered{} \m{7}}
Solve the following puzzle by assigning numeral (0-9) in such a way that each letter is
assigned unique digit which satisfy the following addition. ONE $+$ ONE $+$ TWO $=$ FOUR
\end{asked}

\sbsr{0.48}{%
\lead{The chain}
\begin{itemize}
  \item The addends have three digits and FOUR has four, so \textbf{F is nothing but the
        carry out of column 3}, and $F \geq 1$.
  \item Column 3 is $2O + T + c_2 = 10F + O$, so
  \[ \boxed{\;O + T + c_2 = 10F\;} \]
  \item $O + T \leq 9 + 8 = 17$ and $c_2 \leq 2$, so the left side is at most 19.
        \textbf{$F = 2$ is impossible}, giving $\mathbf{F = 1}$ and $O + T + c_2 = 10$.
  \item Column 1: $2E + O = 10c_1 + R$. \quad Column 2: $2N + W + c_1 = 10c_2 + U$.
  \item \mk{That is everything the constraints force.} Seven of the
        eight letters are free; say so before you choose.
\end{itemize}}{%
\lead{A choice that satisfies all of it}
\begin{itemize}
  \item Take $E = 2$, $O = 6$. Column 1: $4 + 6 = 10$, so $\mathbf{R = 0}$, $c_1 = 1$.
  \item $O + T + c_2 = 10$ with $O = 6$ leaves $T + c_2 = 4$. Take $c_2 = 1$,
        $\mathbf{T = 3}$.
  \item Column 2: $2N + W + 1 = 10c_2 + U$. Take $N = 4$, $W = 8$:
        $8 + 8 + 1 = 17$, so $\mathbf{U = 7}$ and $c_2 = 1$ \checkmark\ as assumed.
\end{itemize}
\par\penalty0\vspace{2pt}
\noindent\colorbox{gdL}{\parbox{\dimexpr\linewidth-2\fboxsep}{\textbf{Answer.}
$O{=}6, N{=}4, E{=}2, T{=}3, W{=}8, F{=}1, U{=}7, R{=}0$, all distinct.
\quad $642 + 642 + 386 = 1670$ \checkmark}}}

%-----------------------------------------------------------------------
\creamq{WRONG $+$ WRONG $=$ RIGHT}{\m{5+3}\m{2+6}}
{\color{sub}\footnotesize \tP{2} \yr{\textbf{70 Bh} \m{5+3}, \textbf{77 Ch} \m{2+6}}
\gap$\cdot$\gap 21 solutions, \mk{nothing at all is forced}
\gap$\cdot$\gap \yr{69 Po} sets \emph{the same eight letters with the two words swapped};
work them as two different puzzles, because they are}

\begin{asked}{\textbf{70 Bh} \textperiodcentered{} Q2 \textperiodcentered{} \m{5+3}}
Solve the following crypto-arithmetic problem, where different letters denote different
integers and identical letters denote same integer. WRONG $+$ WRONG $=$ RIGHT. Explain
the steps that you have followed.
\end{asked}

\sbsr{0.48}{%
\lead{The chain}
\begin{itemize}
  \item \mk{Both sides have five digits}, so there is no sixth carry:
        $2 \times \text{WRONG} < 100000$, hence $\text{WRONG} < 50000$ and
        $\mathbf{W \leq 4}$. \textbf{This is the only bound the puzzle gives you}, and
        writing it down is most of the method mark.
  \item Column 5: $2W + c_4 = R$, so \mk{$R = 2W + c_4$} $-$ $R$ is
        pinned the moment $W$ and that carry are chosen.
  \item Column 1: $2G = 10c_1 + T$, so \mk{$T$ is even}.
  \item Columns 2, 3, 4: $2N + c_1 = 10c_2 + H$, \; $2O + c_2 = 10c_3 + G$, \;
        $2R + c_3 = 10c_4 + I$.
  \item \textbf{No letter is forced to one digit.} State that, then choose.
\end{itemize}}{%
\lead{A choice that satisfies all of it}
\begin{itemize}
  \item Take $W = 3$ with $c_4 = 1$, so $\mathbf{R = 7}$.
  \item Column 1: $G = 1$ gives $\mathbf{T = 2}$, $c_1 = 0$.
  \item Column 2: $N = 8$ gives $16$, so $\mathbf{H = 6}$, $c_2 = 1$.
  \item Column 3: $O = 0$ gives $0 + 1 = 1 = G$ \checkmark, $c_3 = 0$.
  \item Column 4: $2(7) + 0 = 14$, so $\mathbf{I = 4}$ and
        \mk{$c_4 = 1$, which is what column 5 assumed} $-$ that
        closing check is the one step candidates skip.
\end{itemize}
\par\penalty0\vspace{2pt}
\noindent\colorbox{gdL}{\parbox{\dimexpr\linewidth-2\fboxsep}{\textbf{Answer.}
$W{=}3, R{=}7, O{=}0, N{=}8, G{=}1, I{=}4, H{=}6, T{=}2$.
\quad $37081 + 37081 = 74162$ \checkmark}}}

%-----------------------------------------------------------------------
\creamq{SWIM $+$ WEAR $=$ RELAX}{\m{1+6}\m{3+6}}
{\color{sub}\footnotesize \tP{2} \yr{\textbf{74 Bh} \m{1+6}, \textbf{81 Ch} \m{3+6}}
\gap$\cdot$\gap 16 solutions, and \mk{the richest chain of the seven}:
two letters fall out completely and a third is pinned to a fourth}

\begin{asked}{\textbf{74 Bh} \textperiodcentered{} Q2 \textperiodcentered{} \m{1+6}}
Define constraint satisfaction problem (CSP). Solve the following crypto-arithmetic
problem, where different letters denote different integers and identical letters denote
same integer. SWIM $+$ WEAR $=$ RELAX.
\end{asked}

\sbsr{0.52}{%
\lead{The chain}
\begin{itemize}
  \item Two four-digit numbers give a five-digit answer, so the leading digit is only a
        carry: $\mathbf{R = 1}$.
  \item Column 2 is $I + A + c_1 = 10c_2 + A$, and the $A$ cancels:
  \[ \boxed{\;I + c_1 = 10c_2\;} \]
  \item \textbf{Both cases, and one of them dies.} If $c_2 = 1$ then $I + c_1 = 10$,
        forcing $I = 9$, $c_1 = 1$. Column 1 is then $M + R = 10 + X$ with $R = 1$, so
        $M = 9 + X$, which needs $X = 0$ and $M = 9$ $-$
        \mk{but $I$ is already 9}. Contradiction.
  \item So $c_2 = 0$ and $\mathbf{I = 0}$, $c_1 = 0$.
  \item Column 1 with $c_1 = 0$: $M + 1 = X$, so \mk{$X = M + 1$}.
        The two move together.
  \item Column 3: $W + E + c_2 = 10c_3 + L$. \quad Column 4: $S + W + c_3 = 10 + E$.
  \item $A$ meets no constraint but AllDiff $-$ worth saying out loud.
\end{itemize}}{%
\lead{A choice that satisfies all of it}
\begin{itemize}
  \item Take $W = 4$, $E = 3$ with $c_3 = 0$: column 4 gives
        $S = 10 + 3 - 4 = \mathbf{9}$.
  \item Column 3: $4 + 3 + 0 = 7$, so $\mathbf{L = 7}$ and $c_3 = 0$ \checkmark.
  \item $A = 2$, free. $M = 5$, so $X = 6$ by $X = M + 1$.
\end{itemize}
\par\penalty0\vspace{2pt}
\noindent\colorbox{gdL}{\parbox{\dimexpr\linewidth-2\fboxsep}{\textbf{Answer.}
$S{=}9, W{=}4, I{=}0, M{=}5, E{=}3, A{=}2, R{=}1, L{=}7, X{=}6$.
\quad $9405 + 4321 = 13726$ \checkmark}}
\begin{itemize}
  \item \mk{Nine letters, so all ten digits but one are used.} A
        duplicate is easy to make here; check the list before you write the sum.
\end{itemize}}

%-----------------------------------------------------------------------
\creamq{RIGHT $+$ RIGHT $=$ WRONG}{\m{7}}
{\color{sub}\footnotesize \tP{1} \yr{69 Po \m{7}} \gap$\cdot$\gap 11 solutions, nothing
forced \gap$\cdot$\gap \mk{the mirror image of WRONG $+$ WRONG $=$ RIGHT}
above, and it does \emph{not} have the same answer $-$ the letters swap roles, so every
column equation changes}

\begin{asked}{69 Po \textperiodcentered{} Q2 \textperiodcentered{} \m{7}}
Solve the following crypto-arithmetic problem, where different letters denote different
integers and identical letters denote same integer. RIGHT $+$ RIGHT $=$ WRONG. Show all
the step of solving through constrain satisfaction problem.
\end{asked}

\sbsr{0.48}{%
\lead{The chain}
\begin{itemize}
  \item Five digits on both sides, so $2 \times \text{RIGHT} < 100000$ and
        $\mathbf{R \leq 4}$.
  \item Column 5: \mk{$W = 2R + c_4$}.
  \item Column 1: $2T = 10c_1 + G$, so \mk{$G$ is even}.
  \item Columns 2, 3, 4: $2H + c_1 = 10c_2 + N$, \; $2G + c_2 = 10c_3 + O$, \;
        $2I + c_3 = 10c_4 + R$.
  \item Nothing is forced to a single digit.
\end{itemize}
{\color{sub}\footnotesize Note how column 3 feeds on $G$, which column 1 already produced.
Fixing $T$ early therefore settles two columns at once, and that is the cheapest place to
start guessing.}}{%
\lead{A choice that satisfies all of it}
\begin{itemize}
  \item Take $R = 4$ with $c_4 = 0$, so $\mathbf{W = 8}$.
  \item Column 1: $T = 5$ gives $10$, so $\mathbf{G = 0}$, $c_1 = 1$.
  \item Column 2: $H = 6$ gives $13$, so $\mathbf{N = 3}$, $c_2 = 1$.
  \item Column 3: $0 + 1 = 1$, so $\mathbf{O = 1}$, $c_3 = 0$.
  \item Column 4: $2I = 10c_4 + R = 4$, so $\mathbf{I = 2}$ and $c_4 = 0$ \checkmark\
        as column 5 assumed.
\end{itemize}
\par\penalty0\vspace{2pt}
\noindent\colorbox{gdL}{\parbox{\dimexpr\linewidth-2\fboxsep}{\textbf{Answer.}
$R{=}4, I{=}2, G{=}0, H{=}6, T{=}5, W{=}8, O{=}1, N{=}3$.
\quad $42065 + 42065 = 84130$ \checkmark}}}

%-----------------------------------------------------------------------
\creamq{TWO $+$ TWO $=$ FOUR}{\m{2+6}}
{\color{sub}\footnotesize \tP{1} \yr{\textbf{78 Ch} \m{2+6}} \gap$\cdot$\gap
\mk{only 7 solutions}, the tightest of the seven, so a marker here is
more likely than elsewhere to have one in mind $-$ print the count anyway}

\begin{asked}{\textbf{78 Ch} \textperiodcentered{} Q2 \textperiodcentered{} \m{2+6}}
What is well defined problem? Solve the following crypto-arithmetic problem by defining
it. TWO $+$ TWO $=$ FOUR
\end{asked}

\sbsr{0.48}{%
\lead{The chain}
\begin{itemize}
  \item Two three-digit numbers reach at most $999 + 999 = 1998$, so the leading digit of
        a four-digit FOUR is $\mathbf{F = 1}$.
  \item Column 3: $2T + c_2 = 10F + O = 10 + O$. With $c_2 \leq 1$ this needs
        $2T \geq 9$, so \mk{$T \geq 5$}.
  \item Column 1: $2O = 10c_1 + R$, so \mk{$R$ is even}.
  \item Column 2: $2W + c_1 = 10c_2 + U$.
  \item $F$ is the only letter forced.
\end{itemize}}{%
\lead{A choice that satisfies all of it}
\begin{itemize}
  \item Take $T = 7$ with $c_2 = 0$: column 3 gives $14 = 10 + O$, so $\mathbf{O = 4}$.
  \item Column 1: $2(4) = 8$, so $\mathbf{R = 8}$, $c_1 = 0$.
  \item Column 2: $2W + 0 = 10(0) + U$, so $U = 2W$. Take $W = 3$, $\mathbf{U = 6}$,
        and $c_2 = 0$ \checkmark\ as assumed.
\end{itemize}
\par\penalty0\vspace{2pt}
\noindent\colorbox{gdL}{\parbox{\dimexpr\linewidth-2\fboxsep}{\textbf{Answer.}
$T{=}7, W{=}3, O{=}4, F{=}1, U{=}6, R{=}8$.
\quad $734 + 734 = 1468$ \checkmark}}
\begin{itemize}
  \item \mk{Six letters, so this one is quick} $-$ spend the time
        saved on the 2-mark \emph{well-defined problem} half, which is \textbf{Ch 2
        \S2.2} and is where the easy marks are.
\end{itemize}}

%-----------------------------------------------------------------------
\creamq{AB $+$ CD $=$ AAA: what values can B take?}{\m{4}}
{\color{sub}\footnotesize \yr{\textbf{\textit{76 Bh}} \m{4}} \gap$\cdot$\gap
\mk{not a solve-it question}: it asks what is \emph{possible}, so the answer
is a set and a justification}

\begin{asked}{\textbf{\textit{76 Bh}} \textperiodcentered{} \m{4}}
In a correctly worked out crypt-arithmetic problem AB $+$ CD $=$ AAA, what could be the
possible values of B? Justify your answer.
\end{asked}

\sbs{%
\lead{Reasoning}
\begin{itemize}
  \item \textbf{$A = 1$.} Two 2-digit numbers sum to at most $99 + 99 = 198$. The only
        3-digit repdigit AAA in range is \textbf{111}, so $A = 1$.
  \item So $\text{AB} = 10 + B$ and $\text{CD} = 111 - (10 + B) = \mathbf{101 - B}$.
  \item \textbf{$B \geq 3$.} CD must be a 2-digit number, so $101 - B \leq 99$, giving
        $B \geq 2$. And $B = 2$ makes $\text{CD} = 99$, i.e.\ $C = D = 9$, which the
        AllDiff rule forbids. So $B \geq 3$.
  \item \textbf{$B \leq 8$.} For $B$ from 3 to 9, $\text{CD} = 101 - B$ runs
        $98, 97, \dots, 92$, so $C = 9$ throughout. $B = 9$ would make $B = C = 9$.
        \mk{This is the case that is easy to miss}, and it is the whole
        difficulty of the question.
  \item $B = 1$ is out because $A = 1$, and $B = 0$ would need $\text{CD} = 101$.
\end{itemize}}{%
\lead{\mk{Answer: $B \in \{3, 4, 5, 6, 7, 8\}$}, six values}
\vspace{-1pt}
\begin{tabularx}{\linewidth}{@{}L{0.9cm}L{1.2cm}L{1.2cm}X@{}}
\toprule
\hrow \thead{B} & \thead{AB} & \thead{CD} & \thead{Check} \\
\midrule
3 & 13 & 98 & $13 + 98 = 111$ \checkmark \\
4 & 14 & 97 & $14 + 97 = 111$ \checkmark \\
5 & 15 & 96 & $15 + 96 = 111$ \checkmark \\
6 & 16 & 95 & $16 + 95 = 111$ \checkmark \\
7 & 17 & 94 & $17 + 94 = 111$ \checkmark \\
8 & 18 & 93 & $18 + 93 = 111$ \checkmark \\
\midrule
2 & 12 & 99 & \xmark\ $C = D = 9$ \\
9 & 19 & 92 & \xmark\ $B = C = 9$ \\
\bottomrule
\end{tabularx}
\vspace{2pt}
\begin{itemize}
  \item In every case $A = 1$ and $C = 9$; only $B$ and $D$ move, and they move together
        since $D = 8 - (B - 3)$.
\end{itemize}}

%=======================================================================
\T{2.N.3 The water jug problem}

\sbsr{0.55}{%
\lead{The eight production rules, state $(x, y)$}
{\color{sub}\footnotesize $x$ is the contents of the larger jug, $y$ of the smaller.
Written here for the 4-gallon and 3-gallon pair; swap the capacities for any other
version. This is the rule set \emph{Insights} and the lecture decks both use, and it is
what ``write down the production rules'' means.}
\vspace{2pt}
\begin{tabularx}{\linewidth}{@{}L{0.7cm}L{3.5cm}L{2.0cm}X@{}}
\toprule
\hrow \thead{\#} & \thead{Condition} & \thead{Result} & \thead{Meaning} \\
\midrule
1 & $(x,y)$ and $x < 4$ & $(4, y)$ & Fill the 4-gallon jug \\
2 & $(x,y)$ and $y < 3$ & $(x, 3)$ & Fill the 3-gallon jug \\
3 & $(x,y)$ and $x > 0$ & $(0, y)$ & Empty the 4-gallon jug on the ground \\
4 & $(x,y)$ and $y > 0$ & $(x, 0)$ & Empty the 3-gallon jug on the ground \\
5 & $(x,y)$ and $x + y \geq 4$ and $y > 0$ & $(4,\ y - (4 - x))$ & Pour from the 3 into the 4 until the 4 is full \\
6 & $(x,y)$ and $x + y \geq 3$ and $x > 0$ & $(x - (3 - y),\ 3)$ & Pour from the 4 into the 3 until the 3 is full \\
7 & $(x,y)$ and $x + y \leq 4$ and $y > 0$ & $(x + y,\ 0)$ & Pour all of the 3 into the 4 \\
8 & $(x,y)$ and $x + y \leq 3$ and $x > 0$ & $(0,\ x + y)$ & Pour all of the 4 into the 3 \\
\bottomrule
\end{tabularx}}{%
\lead{Formulating it as a state space}
\begin{itemize}
  \item \textbf{State}: the ordered pair $(x, y)$, $0 \leq x \leq 4$, $0 \leq y \leq 3$.
        \textbf{20 states} in all.
  \item \textbf{Initial state}: $(0, 0)$, both jugs empty.
  \item \textbf{Goal state}: $(2, y)$ for any $y$ $-$ the question asks only for 2 gallons
        \emph{in the 4-gallon jug}, so $y$ is unconstrained.
  \item \textbf{Operators}: the eight rules alongside.
  \item \textbf{Path cost}: 1 per rule application.
  \item \mk{Why the state is just a pair}: the jugs have no markings, so
        nothing about \emph{how} the water got there can matter. This is the whole point of
        the exercise and it is worth one sentence in the answer.
  \item \textbf{Problem type}: \textbf{ignorable} (\S2.3) $-$ a wrong move can simply be
        emptied out, so no backtracking machinery is needed.
\end{itemize}}

%-----------------------------------------------------------------------
\creamq{2 gallons in the 4-gallon jug, from a 4 and a 3}{\m{4}\m{2+2+4}}
{\color{sub}\footnotesize \tP{2} \yr{\textbf{\textit{77 Ch}} \m{2+2+4},
\textbf{\textit{71 Bh}} \m{4}} \gap$\cdot$\gap the classic; \emph{Insights} \S2.2.1
works this one}

\begin{asked}{\textbf{\textit{77 Ch}} \textperiodcentered{} \m{2+2+4}}
Assume you are given two jugs, a 4-gallon one and a 3-gallon one, and a pump which has
unlimited water. How can you get exactly 2 gallons of water in the 4-gallon jug? Formalize
the problem, write down Production Rules and draw the Search Tree for the water jug
problem.
\end{asked}

\sbs{%
\lead{Solution, following \emph{Insights}: rules 2, 7, 2, 5, 3, 7}
\vspace{-1pt}
\begin{tabularx}{\linewidth}{@{}L{1.1cm}L{1.5cm}L{0.9cm}X@{}}
\toprule
\hrow \thead{Step} & \thead{State $(x,y)$} & \thead{Rule} & \thead{Action} \\
\midrule
start & $(0, 0)$ & $-$ & both jugs empty \\
1 & $(0, 3)$ & 2 & fill the 3-gallon jug \\
2 & $(3, 0)$ & 7 & pour all of it into the 4-gallon jug \\
3 & $(3, 3)$ & 2 & fill the 3-gallon jug again \\
4 & $(4, 2)$ & 5 & pour into the 4 until it is full; the 4 takes 1, leaving \textbf{2} \\
5 & $(0, 2)$ & 3 & empty the 4-gallon jug \\
6 & $\mathbf{(2, 0)}$ & 7 & pour the 2 gallons into the 4-gallon jug \\
\bottomrule
\end{tabularx}
\vspace{2pt}
\begin{itemize}
  \item \textbf{6 rule applications}, and breadth-first search confirms no shorter path
        exists.
\end{itemize}}{%
\lead{Search tree}
\begin{itemize}
  \item From $(0,0)$ only two rules apply, 1 and 2, so the tree has \textbf{two branches
        at the root}. The 6-step path above is the left branch; the right branch (fill the
        4 first) reaches $(2,3)$ in 6 steps too.
  \item \mk{Draw it as this, not as a bush} $-$ show the two root
        branches and expand the one that succeeds:
\end{itemize}
\vspace{2pt}
{\footnotesize
$(0,0)$ \\
\hspace*{0.6em}$\rightarrow$ [rule 1] $(4,0) \rightarrow (1,3) \rightarrow (1,0)
\rightarrow (0,1) \rightarrow (4,1) \rightarrow \mathbf{(2,3)}$ \\
\hspace*{0.6em}$\rightarrow$ [rule 2] $(0,3) \rightarrow (3,0) \rightarrow (3,3)
\rightarrow (4,2) \rightarrow (0,2) \rightarrow \mathbf{(2,0)}$
}
\vspace{2pt}
\begin{itemize}
  \item \mk{Both are correct and both are 6 moves.} Either scores full
        marks; say which you are drawing.
  \item {\color{sub}\footnotesize \emph{Insights} draws this tree with a self-loop child
        at each level, labelled $(x,y) \rightarrow (x{+}0, y{+}0)$. That is not a legal
        successor $-$ no rule leaves the state unchanged $-$ and it should be left out.}
\end{itemize}}

%-----------------------------------------------------------------------
\creamq{7 gallons in the 9-gallon jug, from a 3, a 5 and a 9}{\m{2+1+2+3}}
{\color{sub}\footnotesize \tP{1} \yr{\textbf{\texttt{79 Bh}} \m{2+1+2+3}} \gap$\cdot$\gap
\mk{the hard variant}: three jugs, and the pump may fill only the smallest}

\begin{asked}{\textbf{\texttt{79 Bh}} \textperiodcentered{} \m{2+1+2+3}}
What are the steps of Problem solving? You are given three empty jugs: a 3-gallon, a
5-gallon and a 9-gallon, and a pump can fill water only in the 3-gallon jug. How do you get
exactly 7 gallons in the 9-gallon jug? Formalize the problem, write production rules and
draw the search tree.
\end{asked}

\sbs{%
\lead{What makes this one different}
\begin{itemize}
  \item \mk{Only the 3-gallon jug may be filled from the pump.} So rule 1
        applies to jug A alone; there is no ``fill the 5'' and no ``fill the 9''.
  \item Every gallon that reaches the 9-gallon jug arrives in multiples of 3, \emph{unless}
        the 5-gallon jug is used to strand a remainder. $9 = 3 \times 3$, so a direct
        route gives only 3, 6 or 9 gallons: \textbf{the 5-gallon jug is not optional, it
        is the whole trick}.
  \item \textbf{State}: the triple $(a, b, c)$ for the 3, 5 and 9-gallon jugs.
        \textbf{Initial} $(0,0,0)$, \textbf{goal} $(a, b, 7)$.
  \item \textbf{Operators}: fill A from the pump; empty any jug; pour $i$ into $j$ until
        $i$ is empty or $j$ is full.
\end{itemize}}{%
\lead{Solution, 9 moves (shortest)}
\vspace{-1pt}
\begin{tabularx}{\linewidth}{@{}L{0.7cm}L{2.1cm}X@{}}
\toprule
\hrow \thead{\#} & \thead{$(a,b,c)$} & \thead{Action} \\
\midrule
 & $(0,0,0)$ & start \\
1 & $(3,0,0)$ & fill A from the pump \\
2 & $(0,3,0)$ & pour A into B \\
3 & $(3,3,0)$ & fill A \\
4 & $(1,5,0)$ & pour A into B; B takes 2 and is full, \mk{leaving 1 in A} \\
5 & $(0,5,1)$ & pour A into C \\
6 & $(3,5,1)$ & fill A \\
7 & $(0,5,4)$ & pour A into C \\
8 & $(3,5,4)$ & fill A \\
9 & $(0,5,\mathbf{7})$ & pour A into C \\
\bottomrule
\end{tabularx}
\vspace{2pt}
\begin{itemize}
  \item \mk{The whole idea is step 4}: fill B to strand exactly 1 gallon in
        A, tip that 1 into C, then add $3 + 3$ on top. $1 + 3 + 3 = 7$.
  \item B is never emptied after step 4; it has done its job.
\end{itemize}}

%-----------------------------------------------------------------------
\creamq{1 litre, from a 5-litre and a 2-litre jug}{\m{8}}
{\color{sub}\footnotesize \tP{1} \yr{81 Ash \m{8}} \gap$\cdot$\gap
this paper asks \mk{three things at once}: is it well defined, list the
attributes for designing an intelligent agent, and solve it. Answer all three}

\begin{asked}{81 Ash \textperiodcentered{} \m{8}}
You are given two jugs, a 5-litre one and a 2-litre one. Neither has any measuring markers
on it. There is a pump that can be used to fill the jugs with water. You are required to
measure exactly 1 litre of water. Is this a well-defined problem? Justify your argument and
list all the attributes to design an intelligent agent. Solve this problem by specifying
states and the operations used.
\end{asked}

\sbsr{0.46}{%
\lead{Part 1: is it well defined? \mk{Yes}}
\begin{itemize}
  \item \textbf{Initial state} given: both jugs empty, $(0,0)$.
  \item \textbf{Operators} given: fill from the pump, empty on the ground, pour one into
        the other.
  \item \textbf{Goal test} given and mechanical: does either jug hold exactly 1 litre?
  \item \textbf{Path cost}: 1 per operation.
  \item Every clause of \S2.2 is satisfied, so it is well defined. It is also
        \textbf{solvable}, since $\gcd(5,2) = 1$ divides 1.
\end{itemize}
\lead{Part 2: attributes for designing an intelligent agent}
\begin{itemize}
  \item This is asking for \textbf{PEAS} plus the environment properties (Ch\,1):
  \begin{itemize}
    \item \textbf{P}: measure exactly 1 litre, in the fewest operations, spilling least.
    \item \textbf{E}: the two jugs, the pump, the ground.
    \item \textbf{A}: fill, empty, pour.
    \item \textbf{S}: whether a jug is full, empty, or neither.
  \end{itemize}
  \item Environment: \textbf{fully observable, deterministic, sequential, static, discrete,
        single agent.} A \textbf{goal-based agent} is the right type, because the answer
        needs a \emph{sequence} planned ahead, not a reflex.
\end{itemize}}{%
\lead{Part 3: solution, 4 moves (shortest)}
\vspace{-1pt}
\begin{tabularx}{\linewidth}{@{}L{0.7cm}L{1.7cm}L{0.8cm}X@{}}
\toprule
\hrow \thead{\#} & \thead{$(x,y)$} & \thead{Rule} & \thead{Action} \\
\midrule
 & $(0,0)$ & $-$ & start \\
1 & $(5,0)$ & 1 & fill the 5-litre jug \\
2 & $(3,2)$ & 6 & pour into the 2-litre jug until full \\
3 & $(3,0)$ & 4 & empty the 2-litre jug \\
4 & $\mathbf{(1,2)}$ & 6 & pour again; the 2 takes 2 more, \mk{leaving 1 in the 5} \\
\bottomrule
\end{tabularx}
\vspace{2pt}
\begin{itemize}
  \item \textbf{The 1 litre ends up in the 5-litre jug}, which is what the question asks
        for: it says ``measure exactly 1 litre'', not ``in the 2-litre jug''.
  \item State space: $6 \times 3 = 18$ states.
  \item \mk{The general rule worth one extra mark}: with jugs of capacity
        $m$ and $n$, exactly the multiples of $\gcd(m,n)$ up to $\max(m,n)$ are
        measurable. Here $\gcd(5,2) = 1$, so \emph{every} whole number of litres from 0 to
        5 can be measured.
\end{itemize}}

%=======================================================================
\T{2.N.4 Farmer, wolf, goat and cabbage}

\creamq{Formulate the river crossing as search and solve it}{\m{8}}
{\color{sub}\footnotesize \tP{1} \yr{73 Ma \m{8}} \gap$\cdot$\gap
the standard example of a problem whose interest is entirely in the
\mk{safety precondition}}

\begin{asked}{73 Ma \textperiodcentered{} \m{8}}
A farmer has a goat, a wolf and a cabbage on the west side of a river and wants to move
everything to the east side. The boat holds the farmer and one other thing. The wolf eats
the goat and the goat eats the cabbage if left alone. Formulate this puzzle as search and
solve it (any method); draw the search tree and final solution.
\end{asked}

\sbs{%
\lead{Formulation}
\begin{itemize}
  \item \textbf{State}: the 4-tuple $(F, W, G, C)$ giving which bank the Farmer, Wolf,
        Goat and Cabbage are on. Write $0$ for west, $1$ for east.
  \item \textbf{Initial state} $(0,0,0,0)$; \textbf{goal state} $(1,1,1,1)$.
  \item \textbf{Operators}: the farmer crosses, taking nothing or exactly one item that is
        currently on his own bank. \textbf{Four operators.}
  \item \textbf{Safety constraint} (this is what makes it a puzzle): a state is illegal if
  \begin{itemize}
    \item the goat and the wolf are together without the farmer, or
    \item the goat and the cabbage are together without the farmer.
  \end{itemize}
  \item \textbf{Path cost}: 1 per crossing.
  \item \textbf{State space}: $2^4 = 16$ states in total, of which
        \mk{only 10 are safe and reachable}. Six are pruned by the
        constraint, and that pruning is the answer to ``how does a constraint reduce the
        search''.
\end{itemize}}{%
\lead{Solution: 7 crossings}
\vspace{-1pt}
\begin{tabularx}{\linewidth}{@{}L{0.7cm}L{3.2cm}X@{}}
\toprule
\hrow \thead{\#} & \thead{$(F,W,G,C)$} & \thead{Crossing} \\
\midrule
 & W W W W & start, everything on the west \\
1 & E W E W & take the \textbf{goat} across \\
2 & W W E W & return \textbf{alone} \\
3 & E E E W & take the \textbf{wolf} across \\
4 & W E W W & \mk{bring the goat back} \\
5 & E E W E & take the \textbf{cabbage} across \\
6 & W E W E & return \textbf{alone} \\
7 & \textbf{E E E E} & take the \textbf{goat} across \\
\bottomrule
\end{tabularx}
\vspace{2pt}
\begin{itemize}
  \item \mk{Step 4 is the step everyone misses.} Nothing else works: after
        the wolf is delivered, the goat cannot stay with it, so it must come back and wait
        while the cabbage goes over.
  \item The goat is the only item that is dangerous \emph{in both directions}, which is
        why it moves first, and why it is the one that gets brought back.
  \item \textbf{7 is the minimum.} There is a mirror-image solution that takes the cabbage
        third instead of the wolf; it is also 7 crossings and equally correct.
\end{itemize}}

%=======================================================================
\T{2.N.5 Tic-tac-toe as a state space}

\creamq{Describe tic-tac-toe in terms of state space, data structure and goal}{\m{7}}
{\color{sub}\footnotesize \tP{1} \yr{\textit{72 Ma} \m{7}} \gap$\cdot$\gap
the paper asks for \mk{complete state space, efficient data structures and
goal}, i.e.\ three named parts. Answer in three parts}

\begin{asked}{\textit{72 Ma} \textperiodcentered{} \m{7}}
Consider a Tic-Tac-Toe game playing as a problem. Prepare a brief description of your
problem in terms of complete state-space, efficient data structures and goal, which are to
be used for a machine to play that game.
\end{asked}

\sbs{%
\lead{1. Complete state space}
\begin{itemize}
  \item A state is an assignment of \textbf{blank}, \textbf{X} or \textbf{O} to each of
        the 9 squares, plus whose turn it is.
  \item \textbf{Naive count}: $3^9 = \mathbf{19{,}683}$ boards. \mk{Most are
        not reachable} $-$ they have five O's and no X's, or two completed lines.
  \item \textbf{Reachable, legal count}: $\mathbf{5{,}478}$ states, counting every board
        that can arise in a real game with X moving first and play stopping at a win.
  \item \textbf{Branching factor} 9 at the root, falling by one each ply; at most
        $9! = 362{,}880$ complete play sequences, which is small enough to search
        \textbf{exhaustively}. That is the point worth making: tic-tac-toe is one of the
        few games where minimax can see to the end of the game.
  \item \textbf{Initial state}: the empty board. \textbf{Operators}: place your mark on
        any blank square. \textbf{Problem type}: \textbf{irrecoverable} (\S2.3) $-$ a
        placed mark cannot be lifted, so the machine must plan before moving.
\end{itemize}}{%
\lead{2. Efficient data structure}
\begin{itemize}
  \item \textbf{A 9-element vector} indexed 1--9, holding 0 for blank, 1 for X, 2 for O.
        Simple, and a move is one array write.
  \item \mk{The efficient version}: two \textbf{9-bit masks}, one for X and
        one for O. Then
  \begin{itemize}
    \item a legal-move test is a bitwise AND against the union,
    \item a win test is \code{mask \& line == line} over the \textbf{8 winning lines}
          (3 rows, 3 columns, 2 diagonals), which is 8 AND operations and no loops,
    \item the whole board fits in one integer, so states can be hashed and cached.
  \end{itemize}
  \item A move-indexed \textbf{lookup table} of best replies is the fastest of all, but it
        is the table-driven agent of Ch\,1: correct, and impractical to build by hand.
\end{itemize}
\lead{3. Goal}
\begin{itemize}
  \item \textbf{Goal test}: any of the 8 lines is filled with one player's mark, or all 9
        squares are taken (a draw).
  \item \textbf{Utility}: $+1$ win, $0$ draw, $-1$ loss, which is what minimax (\S3.4)
        needs.
  \item \mk{The honest closing line}: with both players optimal, tic-tac-toe
        is a \textbf{forced draw}. The machine's real goal is therefore \emph{never to
        lose}, and to win only when the opponent errs.
\end{itemize}}
